\section{The pixel site and the \v{C}ech bridge}

This section records claim nodes and dependency edges.

\begin{proposition}[The pixel-grid site]\label{thm:pixel-site}\lean{Obsmetrologia.Site.PixelGrid.grid_site_grothendieck}\notready\uses{thm:bitopological-site-grothendieck}
\emph{(Open frontier --- stated, not formalized.)} A Grothendieck site on a finite pixel
grid, in the pattern of \cref{thm:bitopological-site-grothendieck}: objects are grid
regions, covers are subdivisions, the three basis axioms discharged by exhaustion. The
pattern transfers; the code does not (the bitopological site's clopen algebra is
hardcoded to the bilattice's own two-point dual), so this is templated new work. This is
the cell reading of a pixel --- Rosenfeld (1970), Khalimsky--Kopperman--Meyer (1990),
Kovalevsky --- given the library's site carrier.
\end{proposition}

\begin{theorem}[Pixel \v{C}ech gluing]\label{thm:pixel-cech-glue}\lean{Obsmetrologia.Cohomology.PixelCech.h1_vanishes}\notready\uses{thm:pixel-site, thm:cech1, thm:cocycle, thm:coboundary, thm:glue, thm:h1-vanishes, thm:sheaf-separated}
\emph{(Open frontier --- stated, not formalized.)} The
\cref{thm:cech1}/\cref{thm:glue}/\cref{thm:h1-vanishes} machinery generalized off
\texttt{PredicationCell} to an abstract cell index and instantiated at the pixel site:
locally agreeing pixel patches (overlapping windows, mipmap levels) glue to a unique
global image, existence by first-covering-patch readout, uniqueness by
\cref{thm:sheaf-separated}. The generalization is mechanical (the existing proofs use
only \texttt{List.find?}, the cocycle hypothesis, and separation) but is real
refactoring, not free reuse.
\end{theorem}

\begin{theorem}[The fibration--\v{C}ech bridge]\label{thm:fibration-cech-bridge}\lean{Obsmetrologia.Cohomology.PixelCech.fibration_cover_glue}\notready\uses{thm:pixelize-adjoint-triple, thm:pixel-cech-glue}
\emph{(Open frontier --- stated, not formalized.)} The technical core of the chapter:
fibration-pixelization is the special case of \v{C}ech gluing whose cover is exactly the
fibers of one base map. The point-sample reading (\cref{thm:pixelize-adjoint-triple},
one signal, one resolution, lossless iff fixed) and the cell reading
(\cref{thm:pixel-cech-glue}, many local patches, global assembly) answer different
questions; this node is the precise statement of how they relate, against the
literature's habitual conflation of the two.
\end{theorem}
